\documentclass[a4paper]{article}     
\usepackage{amsfonts}
\usepackage{amsmath}
\usepackage{amssymb}
\usepackage{graphicx}
\usepackage{color}

\newcommand{\Q}{\mathbb{Q}}
\newcommand{\R}{\mathbb{R}}
\newcommand{\llim}{\lim\limits}
\newcommand{\dint}{\displaystyle\int}

\begin{document}

\noindent \large Auteur: Ilyas Sefiani \\
\noindent \large Studierichting: Informatica\\
\noindent \large Oefeningenreeks 4A nr. 57.2\\

\medskip

\normalsize

$\displaystyle\int \dfrac{dx}{4\cos^2(x)+9\sin^2(x)}$\\

$$\sin^2(x)=\dfrac{1-\cos(2x)}{2}$$ $$\cos^2(x)=\dfrac{1+\cos(2x)}{2}$$\\

$\displaystyle\int \dfrac{2dx}{4(1+\cos(2x))+9(1-\cos(2x))}$\\

$=\displaystyle\int \dfrac{2dx}{13-5\cos(2x)}$\\

$$t=2x$$ $$dt=2dx$$\\

$\displaystyle\int \dfrac{dt}{13-5\cos(t)}$\\

$$u=\tan\left(\dfrac{t}{2}\right)$$ $$dt=\dfrac{2du}{1+u^2}$$ $$\cos(t)=\dfrac{1-u^2}{1+u^2}$$\\

$\displaystyle\int \dfrac{\dfrac{2du}{1+u^2}}{13-5\left(\dfrac{1-u^2}{1+u^2}\right)}$\\

$=\displaystyle\int \dfrac{2du}{8+18u^2}$\\

$=\displaystyle\int \dfrac{du}{4+9u^2}$\\

$$k=3u$$ $$dk=\dfrac{1}{3}dk$$\\

$\dfrac{1}{3}\displaystyle\int \dfrac{du}{4+k^2}$\\

$=\dfrac{1}{6}\operatorname{Bgtan}{\left(\dfrac{k}{2}\right)} + C$\\

$=\dfrac{1}{6}\operatorname{Bgtan}{\left(\dfrac{3u}{2}\right)} + C$\\

$=\dfrac{1}{6}\operatorname{Bgtan}{\left(\dfrac{3\tan\left(\dfrac{t}{2}\right)}{2}\right)} + C$\\

$=\dfrac{1}{6}\operatorname{Bgtan}{\left(\dfrac{3\tan\left(\dfrac{2x}{2}\right)}{2}\right)} + C$\\

$=\dfrac{1}{6}\operatorname{Bgtan}{\left(\dfrac{3\tan(x)}{2}\right)} + C$\\


\begin{thebibliography}{999}
\end{thebibliography}
\end{document} 